% Drawing 23/04-4 -- folder 23, batch 04, archive page 80 (PDF page 81).
% Transcription: à droite, trois croquis : en haut, une boîte en perspective faite de feuillets empilés, avec un point au milieu de sa face avant ; en bas à gauche, un cercle pointé ; en bas à droite, une seconde figure en perspective : une arête oblique dont chaque extrémité porte un éventail de traits, et un côté droit épaissi.
% Kind: two boxes in perspective and a pointed circle. Method 2 (schematic redraw from the facsimile). Status: ok.
% Notes: Three free-hand ink sketches, redrawn stroke by stroke. Top: a box in perspective: its left edge, its bottom edge, four nearly parallel strokes across its front (the top edge and three lines below it, like stacked sheets), a short top back edge, two verticals at the right and a short oblique stroke at the upper right corner; a heavy point in the middle of the front. Lower left: a circle with a point at its centre. Lower right: an oblique edge between two vertices; from the upper vertex strokes to the upper left, to the left, to the upper right (thin) and to the right down to a heavy right side; from the lower vertex strokes to the left, to the lower left (heavy), to the lower right (joining the foot of the heavy side) and a thin one to the upper right. Left out: the long pencil line that runs obliquely across the sheet (a separator), and a faint mark beside the upper vertex of the lower figure. Still differs: what solid each sketch shows is not clear (the note of the transcription says « deux boîtes en perspective »); the strokes are drawn as read, without completing them into a solid. The pointed circle is not mentioned by the note. Overlaid 2026-11-12 (Opus 5.5) on the 400 dpi colour crop: the strokes of the two sketches fall on the page within one or two units; the pointed circle moved 0.7 unit left and made slightly smaller. Redrawn 2026-11-12 (issue 128): to check again. 2026-11-13, Claude Opus 5.5, overlay check (issue 147): laid over the scan rendered at 350 dpi in colour, zoomed on each sketch (13-23 px a unit); not compiled (no TeX engine in the session), the TikZ was rasterised by a script. The box was redrawn stroke by stroke through points read on its ink (its second line ran 2 to 3 units above its ink): the top edge runs into the oblique stroke of the upper right corner; four lines cross the front (the second ends at that oblique stroke, the third turns down at the right into the lower right edge), with a fifth, shorter curve under the third ending in a small hook (not drawn before); the left side is five separate strokes; three verticals at the right; the heavy point and a small dot below it. The pointed circle is an oval about 0.5 unit larger and lower than drawn before, now on its ink. In the lower right figure, the thin stroke from the lower vertex runs on, faint, up to the upper right end of the thin stroke from the upper vertex: added. Everything now lies on its ink within about 1 unit. Left out as before: the long pencil separator line and the faint mark beside the upper vertex of the lower figure. What solids the sketches show remains a reading; the strokes are drawn as they are.
% Coordinates: 1 unit = 1 mm = 1/100 of the crop width (bbox in page fractions: [0.66, 0.25, 0.97, 0.52]).
% Redrawn by Claude (claude-fable-5-2), issue #128, 2026-11-12; not checked by a human; not compiled (no TeX in the session).
% Overlaid by Claude (claude-opus-5-5) 2026-11-12 (issue #128): rasterised from its own coordinates over the 400 dpi colour facsimile and corrected from it. Not compiled. Status: partial.
% Rights: reproduces part of the work -- behind RENDER_TRANSCRIPTS, never published by itself.
\documentclass[tikz,border=3pt]{standalone}
\usepackage{amsmath,amssymb}
\usetikzlibrary{patterns,arrows.meta}
\tikzset{pencil/.style={black!45}, struck/.style={black!35}, mpt/.style={circle,draw,inner sep=1.3pt}, pt/.style={circle,fill,inner sep=0.9pt}, opt/.style={circle,draw,fill=white,inner sep=1.1pt}}
\begin{document}
\begin{tikzpicture}[x=1mm,y=1mm,line cap=round,line join=round,font=\scriptsize]
% the box: top edge running into the oblique stroke at the upper right corner
\draw[thick] plot[smooth] coordinates {(58.4,-11) (66.4,-12.1) (75,-14) (81.5,-16) (84.9,-17.1) (87.7,-16.6)};
\draw (87.7,-16.6) -- (85.8,-19.6) -- (83.2,-24) -- (80,-27.4);
% the stacked sheets across the front
\draw plot[smooth] coordinates {(53.3,-17.7) (54.8,-16.8) (58.2,-16.6) (64.2,-17.9) (70.7,-19.6) (77.2,-21.8) (83.2,-24)};
\draw plot[smooth] coordinates {(52.4,-27.6) (55.6,-26.5) (59.9,-25.2) (65.1,-24.8) (70.7,-26.1) (77.2,-28.3) (83.6,-28.5) (88,-30.2)};
\draw plot[smooth] coordinates {(55.8,-27) (58.2,-26.3) (62.1,-27.8) (66.4,-29.3) (70.7,-30.2) (73.7,-30) (74.9,-29.1)};
\draw plot[smooth] coordinates {(54.3,-36) (56.1,-36.4) (60.8,-35.1) (66.4,-35.8) (72.9,-36.9) (78,-37.3) (84.1,-38.8)};
\draw (88,-30.2) -- (85.8,-33.9) -- (83.6,-38.2) -- (81.7,-43.1);
% bottom edge
\draw plot[smooth] coordinates {(48.9,-43.1) (55.6,-44.6) (64.2,-45.7) (75.4,-46.8)};
% the left side, in five strokes
\draw plot[smooth] coordinates {(55.3,-12.7) (52.6,-19.2) (50.7,-24.6)};
\draw plot[smooth] coordinates {(55.2,-14.9) (54.5,-18.8) (54.1,-24.8)};
\draw plot[smooth] coordinates {(52.8,-18.8) (51.3,-26) (51.1,-34.7)};
\draw (50.7,-27.4) -- (48.9,-36.4);
\draw plot[smooth] coordinates {(53.5,-27) (51.7,-30.4) (50.4,-35.6) (49.6,-40.8) (48.9,-43.1)};
% three verticals at the right
\draw (83.2,-21.6) -- (83,-38.2);
\draw (86,-17.3) -- (86,-28.7);
\draw (78,-31.3) -- (77.9,-46.4);
\node[pt] at (67.7,-32.4) {};
\fill (69.5,-35.2) circle[radius=0.25];
\draw (22.8,-100.9) ellipse[x radius=9.7, y radius=8.5];
\node[pt] at (22.4,-100.5) {};
\draw (65.5,-84.5) -- (62,-100) -- (59.5,-110);
\draw plot[smooth, tension=0.55] coordinates {(65.5,-84.5) (58,-81) (51.5,-79.5)};
\draw (65.5,-84.5) -- (56,-87);
\draw[pencil] (65.5,-84.5) -- (76.5,-79.5);
\draw[pencil] (71.2,-102.4) -- (76.7,-79.4);
\draw plot[smooth, tension=0.55] coordinates {(65.5,-84.5) (70,-87) (74,-92)};
\draw[very thick] (74,-92) -- (72.5,-104) -- (69.5,-115);
\draw (59.5,-110) -- (42,-109.5);
\draw[very thick] (59.5,-110) -- (47,-116.5);
\draw plot[smooth, tension=0.55] coordinates {(59.5,-110) (64,-112) (69.5,-115)};
\draw[pencil] (59.5,-110) -- (71,-102);
\end{tikzpicture}
\end{document}
